\chapter{The retained model}
\label{ch:ArchonFull600Theory_RetainedModel}

This chapter is Milestone 5 of the handoff: the census table O9 as a
checked input, the cut chambers and cap signatures, the retained model
with the hypotheses R, F, I, C and E as inputs, Lemmas 6, 7 and 8 (with N4
and N5 as decidable checks), and Lemmas 9 to 11 on the nine selected
orbits (with N11). The hypotheses are not proved: their truth for the
real puzzle is evidence, not theory. Source line numbers refer to
\texttt{research/Full\_600cell\_Puzzle\_Theory.tex}.

\section{The census table}

\begin{definition}[The census table]
  \label{def:census}
  \lean{ArchonFull600Theory.FrameGroupType, ArchonFull600Theory.CensusRow,
    ArchonFull600Theory.censusRows, ArchonFull600Theory.census,
    ArchonFull600Theory.checkCensusJson}
  \leanok
  % SOURCE: hypothesis R, lines 786--829
  % SOURCE QUOTE: Orbit & Rank & Stickers per piece & Pieces \(N_o\) & Frame group
  % \(H_o\)
  \textit{Source: hypothesis R, the table.}
  The \(35\) rows give the rank, the stickers per piece \(k_o\), the
  pieces \(N_o\) and the frame group type \(H_o\in\{1,C_2,C_5,D_5,A_5\}\).
  They are typed in. Compiling the module reads
  \texttt{assets/census.json} and stops the build unless every row and
  the file's totals agree with it.
\end{definition}

\begin{lemma}[Census sums and frame types]
  \label{lem:census-sums}
  \lean{ArchonFull600Theory.censusSums, ArchonFull600Theory.frameTypeCounts,
    ArchonFull600Theory.invariantOrder, ArchonFull600Theory.FrameGroupType.card_carrier,
    ArchonFull600Theory.FrameGroupType.card_abelianization}
  \leanok
  \uses{def:census, lem:n5-dihedral, lem:n5-alternating}
  % SOURCE: hypothesis R, lines 832--836; Lemma 8, lines 1042--1070
  % SOURCE QUOTE: \sum_oN_o=176520,\qquad 600+\sum_ok_oN_o=259800.
  \textit{Source: hypothesis R; N4 and N5.}
  \(\sum_oN_o=176520\), \(600+\sum_ok_oN_o=259800\) and
  \(600+\sum_oN_o=177120\). Seven orbits have \(C_2\), one \(D_5\), two
  \(C_5\), one \(A_5\) and \(24\) the trivial group. The abelianizations
  have orders \(2,5,2,1,1\) for \(C_2,C_5,D_5,A_5,1\), so
  \(\prod_o2|H_o^{\mathrm{ab}}|=2^{43}5^2\).
\end{lemma}
\begin{proof}
  \leanok
  Kernel evaluation over the table. \(|D_5^{\mathrm{ab}}|=10/5\) and
  \(|A_5^{\mathrm{ab}}|=60/60\) by N5.
\end{proof}

\begin{lemma}[N11: the selected orbits]
  \label{lem:selected-census}
  \lean{ArchonFull600Theory.selectedOrbitCensus}
  \leanok
  \uses{def:census}
  % SOURCE: hypothesis C, lines 862--864
  % SOURCE QUOTE: The solved coloring has \(600\) distinct home-cell symbols, with six
  % copies of each in O01, four in O33, and twelve in each of the other seven selected orbits.
  \textit{Source: hypothesis C and Lemma 10.}
  The orbits \(\{1,4,5,7,18,19,27,31,33\}\) have one sticker per piece and
  \(N_o=600m_o\) pieces, with \(m_o=6\) for O01, \(4\) for O33 and \(12\)
  otherwise.
\end{lemma}
\begin{proof}
  \leanok
  Kernel evaluation over the table.
\end{proof}

\section{Caps, cut chambers and signatures}

\begin{definition}[Caps and cut chambers]
  \label{def:cut-chambers}
  \lean{ArchonFull600Theory.Cap, ArchonFull600Theory.capDepth, ArchonFull600Theory.capPole,
    ArchonFull600Theory.cutValue, ArchonFull600Theory.capRegion,
    ArchonFull600Theory.cutComplement, ArchonFull600Theory.chamber,
    ArchonFull600Theory.completeSignature}
  \leanok
  \uses{lem:facet-pole, def:symmetry-group}
  % SOURCE: retained-model definitions, lines 738--769
  % SOURCE QUOTE: \alpha=\frac{121}{125}, \qquad U_c=\{x\in\widetilde P:n_c\cdot
  % x>\alpha\|n_c\|^2\}.
  \textit{Source: the retained-model definitions.}
  The \(600\) caps are the four-cliques of the edge graph, with poles
  \(n_c\) and \(\alpha=121/125\). The cap region is
  \(U_c=\{x\in P:n_c\cdot x>\alpha\|n_c\|^2\}\). A chamber is a choice of
  strict side of every cut hyperplane intersected with
  \(\operatorname{int}P\), and the complete signature of a point is
  \(M(x)=\{c:x\in U_c\}\). The model here is \(P\); the cell-pole model
  \(\widetilde P\) of Lemma 5 multiplies every cut function by
  \((4/\phi)^2\) (\texttt{cutValue\_cellPoleMap}), so it has the same
  chambers and signatures. \(K\) acts on the caps through their vertices.
\end{definition}

\begin{theorem}[Lemma 6, geometric part]
  \label{thm:signatures}
  \lean{ArchonFull600Theory.convex_chamber, ArchonFull600Theory.completeSignature_eq_iff,
    ArchonFull600Theory.completeSignature_smul, ArchonFull600Theory.isPreconnected_patch,
    ArchonFull600Theory.cutValue_cellPoleMap}
  \leanok
  \uses{def:cut-chambers, thm:cell-pole}
  % SOURCE: retained-model definitions, lines 751--755; Lemma 6, lines 914--935
  % SOURCE QUOTE: A nonempty such intersection is convex and hence connected.
  \textit{Source: the retained-model definitions and Lemma 6.}
  Chambers are convex. Two points of \(\operatorname{int}P\) off every cut
  hyperplane have the same complete signature exactly when they lie in the
  same connected component of that set, so complete signatures identify
  the cut-chamber positions. For \(k\in K\), \(M(kx)=kM(x)\). A chamber
  closure meets a closed facet in a convex set. The source's patch, this
  set with its lower-dimensional boundary removed, is not stated.
\end{theorem}
\begin{proof}
  \leanok
  Cut functions are affine, so the chamber conditions are convex. Along a
  preconnected set off the cuts no cut function changes sign, by the
  intermediate value theorem, so the signature is constant on components;
  conversely equal signatures put both points in one convex chamber. \(K\)
  preserves inner products and permutes the poles with the caps.
\end{proof}

\section{The retained model and its hypotheses}

\begin{definition}[The retained model and hypothesis R]
  \label{def:retained-model}
  \lean{ArchonFull600Theory.RetainedModel, ArchonFull600Theory.RetainedModel.legalGroup,
    ArchonFull600Theory.HypothesisR, ArchonFull600Theory.HypothesisR.orbitStickers}
  \leanok
  \uses{def:census, def:cut-chambers}
  % SOURCE: hypothesis R, lines 776--784
  % SOURCE QUOTE: The retained model is the complete surface subdivision just described, with
  % no hidden identification of distinct pieces having the same recorded address.
  \textit{Source: hypothesis R.}
  The data are sticker slots, pieces, the sticker-to-piece map \(\rho\) and
  the primitives \(H_c,T_c\) of each cap; \(G=\langle H_c,T_c\rangle\). R
  says: legal moves carry pieces to pieces; \(600\) one-sticker centres
  are fixed by every primitive; the \(35\) moving classes have the census
  sizes, each is one orbit of \(G\) on pieces, and legal moves keep each
  class; each piece has a point of its chamber and distinct pieces are
  distinct chambers; each sticker has a host facet that determines it
  within its piece; ranks are the census ranks; and each orbit has a frame
  group \(H_o\le S_{k_o}\) of the census type containing every legal
  return map on a reference piece.
\end{definition}

\begin{definition}[Hypotheses F, C and E]
  \label{def:hypotheses-fce}
  \lean{ArchonFull600Theory.HypothesisF, ArchonFull600Theory.HypothesisC,
    ArchonFull600Theory.HypothesisE}
  \leanok
  \uses{def:retained-model, def:pure-controller, def:symmetry-group}
  % SOURCE: hypotheses F, C and E, lines 839--872
  % SOURCE QUOTE: Every cap admits its full tetrahedral rotation group \(A_4\), acting on the
  % same complete subdivision.
  \textit{Source: hypotheses F, C and E.}
  F: each cap's \(A_4\), the even permutations of its four vertices, acts on
  the sticker slots, and \(H_c\), \(T_c\) are the images of an element of
  order \(2\) and one of order \(3\). C: the nine selected orbits have
  pure-controller certificates in \(G\), and the solved colouring has
  \(600\) symbols with \(6\), \(4\) and \(12\) copies of each. E: \(K\) acts
  faithfully on the same sticker slots, respects pieces, sends the centre
  of cap \(c\) to the centre of \(k\cdot c\), and conjugates each rotation
  of cap \(c\) into the rotation of cap \(k\cdot c\) that permutes its
  vertices through \(k\).
\end{definition}

\begin{lemma}[The model's totals]
  \label{lem:model-totals}
  \lean{ArchonFull600Theory.HypothesisR.card_piece, ArchonFull600Theory.HypothesisR.card_sticker}
  \leanok
  \uses{def:retained-model, lem:census-sums}
  % SOURCE: hypothesis R, lines 780--781
  % SOURCE QUOTE: It has \(259800\) sticker slots, \(177120\) surface pieces
  \textit{Source: hypothesis R.}
  Under R there are \(177120\) pieces and \(259800\) sticker slots.
\end{lemma}
\begin{proof}
  \leanok
  Count pieces by class and stickers by piece, and use the census sums.
\end{proof}

\begin{theorem}[Lemma 6]
  \label{thm:lemma-six}
  \lean{ArchonFull600Theory.alternatingFour_words, ArchonFull600Theory.closure_halfTurn_thirdTurn,
    ArchonFull600Theory.HypothesisF.closure_primitives,
    ArchonFull600Theory.HypothesisR.pieceAction, ArchonFull600Theory.HypothesisR.piece_legal,
    ArchonFull600Theory.HypothesisR.signature_injective,
    ArchonFull600Theory.HypothesisR.address_injective,
    ArchonFull600Theory.HypothesisR.rank_multiplicity}
  \leanok
  \uses{def:retained-model, def:hypotheses-fce, thm:signatures}
  % SOURCE: Lemma 6, lines 888--903
  % SOURCE QUOTE: Under R and F, \(H_c,T_c\) generate the full \(A_4\) action at each cap.
  \textit{Source: Lemma 6.}
  Under F, \(H_c\) and \(T_c\) generate the cap's \(A_4\) action. Under R,
  \(\rho(g(s))=\bar g(\rho(s))\) for a homomorphism \(g\mapsto\bar g\) from
  \(G\); distinct pieces have distinct complete signatures; \((M(A),h)\)
  addresses the sticker slots; and \(r(A)=|M(A)|-1\) and
  \(k(A)=|\operatorname{Host}(A)|\) are the census rank and \(k_o\).
\end{theorem}
\begin{proof}
  \leanok
  A kernel check over \(A_4\) shows that an element of order \(2\) and one
  of order \(3\) generate it, the source's Klein four-group argument; an
  equivalence moves it to each cap. Legal moves respect pieces, so they
  induce a permutation of pieces. Signature injectivity follows from the
  geometric part and R's distinct chambers, and addresses from it and R's
  host condition.
\end{proof}

\section{Lemma 7: transported frames}

\begin{definition}[The ambient wreath product]
  \label{def:frame-wreath}
  \lean{ArchonFull600Theory.FrameWreath, ArchonFull600Theory.FrameWreath.mul_perm,
    ArchonFull600Theory.FrameWreath.mul_frame, ArchonFull600Theory.FrameWreath.natCard}
  \leanok
  % SOURCE: Lemma 7, lines 962--966
  % SOURCE QUOTE: (p,a)(r,b)=(r\circ p,c), \qquad c_i=a_i b_{p(i)}.
  \textit{Source: Lemma 7.}
  \(H^X\rtimes S_X\) has elements \((p,a)\), and
  \(|H^X\rtimes S_X|=|X|!\,|H|^{|X|}\). Mathlib composes right to left, so
  \((y\,x).p=y.p\circ x.p\) and \((y\,x).a_i=y.a_{x.p(i)}\,x.a_i\): this is
  the source's chronological law with \(x\) applied first.
\end{definition}

\begin{theorem}[Lemma 7]
  \label{thm:lemma-seven}
  \lean{ArchonFull600Theory.HypothesisR.CoherentFrames, ArchonFull600Theory.HypothesisR.coherentFrames,
    ArchonFull600Theory.HypothesisR.CoherentFrames.coord,
    ArchonFull600Theory.HypothesisR.CoherentFrames.coord_apply,
    ArchonFull600Theory.HypothesisR.CoherentFrames.coord_injective,
    ArchonFull600Theory.HypothesisR.transportedFrames,
    ArchonFull600Theory.HypothesisR.natCard_ambient}
  \leanok
  \uses{def:retained-model, def:frame-wreath}
  % SOURCE: Lemma 7, lines 954--961
  % SOURCE QUOTE: Under R, coherent transported frame choices give an injective homomorphism
  % \begin{equation} G\hookrightarrow \prod_{o=0}^{34}\left(H_o^{N_o}\rtimes S_{N_o}\right).
  \textit{Source: Lemma 7.}
  Under R, coherent transported frames exist and give an injective
  homomorphism \(G\hookrightarrow\prod_o(H_o^{N_o}\rtimes S_{N_o})\) with
  \(g(\mathrm{frame}(i,f))=\mathrm{frame}(p(i),a_if)\). The ambient product
  has order \(\prod_oN_o!\,|H_o|^{N_o}\). The remark that a change of
  frames conjugates the coordinates is not stated.
\end{theorem}
\begin{proof}
  \leanok
  Fix a reference piece and frame in each orbit and a legal transporter to
  every position. The increment of \(g\) at \(i\) is the return map
  \(t_{g(i)}^{-1}gt_i\) on the reference piece, which R puts in \(H_o\).
  The coordinate map is a homomorphism by the wreath law and injective
  because frames cover every moving sticker and the centres are fixed.
\end{proof}

\section{Lemma 8: invariants and the upper bound}

\begin{lemma}[Lemma 8: the invariants]
  \label{lem:invariants}
  \lean{ArchonFull600Theory.FrameWreath.orientation, ArchonFull600Theory.FrameWreath.positionSign,
    ArchonFull600Theory.FrameWreath.invariant, ArchonFull600Theory.FrameWreath.invariant_surjective,
    ArchonFull600Theory.invariantPi, ArchonFull600Theory.invariantPi_surjective,
    ArchonFull600Theory.index_ker_invariantPi,
    ArchonFull600Theory.HypothesisR.natCard_invariantTarget,
    ArchonFull600Theory.HypothesisR.index_invariantKernel}
  \leanok
  \uses{def:frame-wreath, lem:census-sums, thm:lemma-seven}
  % SOURCE: Lemma 8, lines 1023--1075
  % SOURCE QUOTE: \Phi_o(p,a)=\sum_i\chi_o(a_i).
  \textit{Source: Lemma 8.}
  On \(H^X\rtimes S_X\), \(\operatorname{sgn}p\) and
  \(\Phi(p,a)=\sum_i\chi(a_i)\in H^{\mathrm{ab}}\) are homomorphisms, and
  together they map onto \(\{\pm1\}\times H^{\mathrm{ab}}\) when
  \(|X|\ge2\). Over the \(35\) orbits the joint invariant maps onto a group
  of order \(2^{43}5^2\), so its kernel has exactly this index.
\end{lemma}
\begin{proof}
  \leanok
  A permutation of \(X\) reindexes the abelian sum. A transposition sets
  the sign and a single frame coordinate sets \(\Phi\). The target order
  is \(\prod_o2|H_o^{\mathrm{ab}}|\), by the census types.
\end{proof}

\begin{theorem}[Lemma 8: the conditional upper bound]
  \label{thm:invariant-bound}
  \lean{ArchonFull600Theory.HypothesisI, ArchonFull600Theory.HypothesisI.invariantPi_coord,
    ArchonFull600Theory.HypothesisI.invariantBound}
  \leanok
  \uses{lem:invariants, thm:lemma-seven}
  % SOURCE: hypothesis I, lines 844--849; Lemma 8, lines 1011--1018
  % SOURCE QUOTE: Under R and I, \begin{equation} |G| \le
  % \frac{\displaystyle\prod_{o=0}^{34}N_o!\,|H_o|^{N_o}} {2^{43}5^2}.
  \textit{Source: hypothesis I and Lemma 8.}
  I: in some coherent transported frames every \(H_c\) and \(T_c\) has even
  positional permutation and trivial orientation image in every orbit.
  Under R and I,
  \(|G|\cdot2^{43}5^2\le\prod_oN_o!\,|H_o|^{N_o}\), so
  \(|G|\le\prod_oN_o!\,|H_o|^{N_o}/(2^{43}5^2)\), and the denominator is
  the exact index of the invariant kernel. Equality is not asserted.
\end{theorem}
\begin{proof}
  \leanok
  By closure induction every legal element lies in the kernel, so \(G\)
  embeds in it, and \(|\ker|\cdot\mathrm{index}=|\text{ambient}|\).
\end{proof}

\section{Lemmas 9 to 11 on the selected orbits}

\begin{lemma}[The selected orbits in the model]
  \label{lem:selected-model}
  \lean{ArchonFull600Theory.HypothesisR.card_orbitStickers,
    ArchonFull600Theory.HypothesisR.card_selectedOrbitStickers,
    ArchonFull600Theory.HypothesisR.selected_single_sticker,
    ArchonFull600Theory.HypothesisR.selectedOrbitStickers_disjoint,
    ArchonFull600Theory.HypothesisR.legal_mem_selectedOrbitStickers}
  \leanok
  \uses{def:retained-model, lem:selected-census}
  % SOURCE: Lemma 9, lines 1087--1090
  % SOURCE QUOTE: on the nine selected single-sticker orbits and fixing every other sticker.
  \textit{Source: Lemma 9 and N11.}
  Under R, \(|O_o|=k_oN_o\); the nine selected orbits have one sticker per
  piece and \(3600\), \(2400\) and \(7200\) sticker slots; they are pairwise
  disjoint and invariant under \(G\).
\end{lemma}
\begin{proof}
  \leanok
  Count the stickers of \(O_o\) piece by piece; the rest is R and the
  census.
\end{proof}

\begin{lemma}[The cap alphabet realizes \(G\)]
  \label{lem:cap-words}
  \lean{ArchonFull600Theory.foldl_moves, ArchonFull600Theory.wordPermutation_append,
    ArchonFull600Theory.exists_word_pow, ArchonFull600Theory.RetainedModel.capH,
    ArchonFull600Theory.RetainedModel.capT, ArchonFull600Theory.RetainedModel.capWord_surjective}
  \leanok
  \uses{def:retained-model, def:word-coloring}
  % SOURCE: Lemma 11, lines 1193--1197
  % SOURCE QUOTE: \mathcal A=\{H_c,T_c,T_c^{-1}:c\in\mathcal C\}, \qquad q=1800.
  \textit{Source: Lemma 11.}
  For any enumeration of the caps, every legal element is the value of a
  word in \(\{H_c,T_c,T_c^{-1}\}\).
\end{lemma}
\begin{proof}
  \leanok
  Words compose by concatenation and include every generator. In the
  finite group \(G\), \(x^{-1}=x^{\operatorname{ord}x-1}\) is a positive
  power.
\end{proof}

\begin{theorem}[Lemmas 9 to 11 under R and C]
  \label{thm:selected-lemmas}
  \lean{ArchonFull600Theory.HypothesisC.supportedAlternating,
    ArchonFull600Theory.HypothesisC.balancedColorBound,
    ArchonFull600Theory.HypothesisC.wordDistanceBound,
    ArchonFull600Theory.HypothesisC.totalVariationBound}
  \leanok
  \uses{lem:selected-model, lem:cap-words, def:hypotheses-fce, lem:supported-alternating-actions,
    lem:balanced-color-lower-bound, lem:word-distance-lower-bound,
    lem:total-variation-lower-bound}
  % SOURCE: Lemma 9, lines 1087--1090; Lemma 10, lines 1138--1147; Lemma 11, lines 1191--1202
  % SOURCE QUOTE: Under R and C, \(G\) contains a subgroup acting as \begin{equation}
  % A_{3600}\times A_{2400}\times A_{7200}^{\,7} \end{equation}
  \textit{Source: Lemmas 9, 10 and 11.}
  Under R and C: \(G\) contains
  \(A_{3600}\times A_{2400}\times A_{7200}^7\) acting on the nine selected
  orbits and fixing every other sticker; there are at least \(B\) reachable
  face-colour states, \(B\) of them with every unselected sticker solved;
  some state is at word distance at least \(46648\); and a law supported
  on the states of \(1000\)-letter words is at total variation more than
  \(1-10^{-148595}\) from the uniform law.
\end{theorem}
\begin{proof}
  \leanok
  Milestone 1's abstract theorems, with the selected orbits, the
  certificates and the colouring from R and C, and the realization of
  \(G\) by words proved above.
\end{proof}
